\documentclass[11pt]{article}

\usepackage[T1]{fontenc}
\usepackage{lmodern}
\usepackage{amsmath,amssymb,amsthm}
\usepackage[margin=1in]{geometry}
\usepackage{microtype}
\usepackage[hidelinks]{hyperref}

\newtheorem{theorem}{Theorem}
\newtheorem{lemma}[theorem]{Lemma}
\newcommand{\Tfour}{T_4}

\title{The Exact Forced-Win Time on the $4\times4$\\
Transversal Achievement Board}
\author{Xiaojian Zeng\\
\small PBC School of Finance, Tsinghua University, Beijing 100190, China\\
\small \texttt{xiaojianzeng@foxmail.com}}
\date{}

\hypersetup{
  pdftitle={The Exact Forced-Win Time on the 4 x 4 Transversal Achievement Board},
  pdfauthor={Xiaojian Zeng},
  pdfsubject={Exact minimax analysis of a finite positional game},
  pdfkeywords={transversal achievement game, positional game, exact certificate, minimax}
}

\begin{document}
\maketitle

\begin{abstract}
Two players alternately claim empty cells of a $4\times4$ board, with X
moving first, and the first player to own four cells in distinct rows and
columns wins.  We prove that X can force a win by ply $11$, but not by ply
$9$.  Thus the minimax forced-win time is $\Tfour=11$.  The upper bound is
Guan's explicit strategy.  For the lower bound we give an exact positional O
policy with $6{,}861$ choices.  A separately implemented verifier reconstructs
the $24$ winning transversals, branches over every legal X move, checks every
certified O reply and every terminal position, and finds no X win through ply
$9$.  The result concerns only the fixed $4\times4$ game; it does not settle
the proposed general formula for all board sizes.
\end{abstract}

\section{Context and statement}

Erickson posed the transversal achievement game for an $n\times n$ grid
\cite{Erickson2010}.  A transversal contains one cell in every row and every
column.  Ranđelović proved that the first player wins for $n\geq4$
\cite{Randjelovic2025}.  Guan subsequently gave a fully specified strategy
that wins by ply $2n+3$ for every $n\geq4$ \cite{Guan2026}.  For $n=4$ this
gives an upper bound of $11$ plies.  Guan's exhaustive search shows that this
bound is attained against that particular strategy for $n=4,5,6$, and
Conjecture~14 of \cite{Guan2026} proposes that $2n+3$ is the optimal-play
length.  As emphasized there, checking the replies to one X strategy does not
exclude a different, faster X strategy.

We determine the complete minimax value for $n=4$.  Let
\[
 B=[4]\times[4],\qquad
 \mathcal T=\bigl\{\{(i,\pi(i)):i\in[4]\}:\pi\in S_4\bigr\}.
\]
Players X and O alternately claim a free cell, X on odd plies and O on even
plies, and play stops as soon as the mover's claimed set contains a member of
$\mathcal T$.  For pure strategies $\sigma_X,\sigma_O$, let
$\tau_X(\sigma_X,\sigma_O)$ be X's winning ply and set it to $\infty$ when O
wins or the game is drawn.  Define
\[
 \Tfour=\min_{\sigma_X}\max_{\sigma_O}
 \tau_X(\sigma_X,\sigma_O).
\]

\begin{theorem}\label{thm:main}
For the $4\times4$ transversal achievement game, $\Tfour=11$.
\end{theorem}

The only contribution claimed here is this fixed-board determination.  In
particular, no conclusion is drawn for Guan's conjecture when $n>4$.

\section{The horizon-nine defense certificate}

Index the cells by $4r+c$, where $0\leq r,c<4$, and encode each player's
claimed cells by a $16$-bit integer.  For $\pi\in S_4$, the corresponding
winning mask is
\[
 W_\pi=\sum_{r=0}^{3}2^{4r+\pi(r)}.
\]
There are exactly $4!=24$ such masks.  A position is a pair $(x,o)$ of
disjoint ownership masks.  Because legal moves, terminal status and future
winning conditions depend only on $(x,o)$, histories that reach the same pair
may be merged.

The ancillary file \texttt{horizon9\_policy.json} has one record
\[
 (\texttt{x\_mask},\texttt{o\_mask},\texttt{move})
\]
for every nonterminal O-to-move position reachable while the policy is
followed.  Its move is a cell index.  The implicit directed graph starts at
$(0,0)$, contains every free-cell successor at each X node, and contains the
single recorded successor at each O node.  Each edge claims one cell, so the
graph is acyclic.

The following elementary certificate lemma isolates the logical content of
the computation.

\begin{lemma}\label{lem:certificate}
Suppose a finite policy graph rooted at the empty position has the following
properties through ply $9$:
\begin{enumerate}
 \item every state has disjoint in-board masks with the ownership counts
 required by alternating play;
 \item wins are tested before a state is treated as a horizon leaf;
 \item every legal X move is represented at every reachable nonterminal X
 node;
 \item every reachable nonterminal O node has one represented legal O move;
 and
 \item no reachable terminal state is an X win.
\end{enumerate}
Then O has a strategy that prevents X from winning through ply $9$.
\end{lemma}

\begin{proof}
Order the graph by decreasing ply.  At a safe ply-$9$ leaf, X has not won.  An
earlier O win also prevents an X win.  At an O node, the represented legal
move leads to a safe successor.  At an X node, every legal move leads to a
safe successor.  Backward induction therefore marks the root safe against
every sequence of X moves.  The recorded choices define O on all positions
reachable under this policy and can be extended arbitrarily on unreachable O
positions, giving a strategy in the usual game-theoretic sense.
\end{proof}

The independent verifier \texttt{verify\_horizon9\_policy.py} checks precisely
these hypotheses without importing the policy generator.  It reconstructs
the $24$ masks $W_\pi$, rejects duplicate or unreachable policy records,
checks mask bounds, disjointness, move legality and turn counts, and branches
over every free cell at X nodes.  It tests both players' wins before accepting
a horizon leaf.  Its exhaustive replay gives the counts in
Table~\ref{tab:counts}.

\begin{table}[ht]
\centering
\caption{Exhaustive replay of the certified O policy.}
\label{tab:counts}
\begin{tabular}{lr}
Quantity & Count\\
Reachable positions & $28{,}604$\\
Exhaustively branched X positions & $4{,}911$\\
Checked legal X edges & $41{,}768$\\
O positions with a checked policy choice & $6{,}861$\\
Checked O edges & $6{,}861$\\
Terminal O-win positions & $111$\\
Safe ply-$9$ positions & $16{,}721$\\
Terminal X-win positions & $0$\\
\end{tabular}
\end{table}

Thus Lemma~\ref{lem:certificate} supplies the universal quantifier absent
from a search tied to a prescribed X strategy: one O policy answers every
legal X branch through the relevant horizon.  No symmetry quotient is used;
all masks and moves refer to the original labelled board.

\begin{proof}[Proof of Theorem~\ref{thm:main}]
The verified policy and Lemma~\ref{lem:certificate} show that X cannot force a
win by ply $9$.  Hence $\Tfour>9$.  X moves only on odd plies, while Guan's
strategy gives $\Tfour\leq11$ \cite{Guan2026}.  Therefore $\Tfour=11$.
\end{proof}

\section{Reproduction, checks and limitations}

From the directory containing the ancillary \texttt{evidence} folder, the
certificate can be regenerated and then checked with
\begin{verbatim}
python3 evidence/horizon9_solver.py \
  evidence/horizon9_policy.json evidence/horizon9_solver_summary.json
python3 evidence/verify_horizon9_policy.py \
  evidence/horizon9_policy.json evidence/horizon9_verification.json
\end{verbatim}
The generator is an exact unsymmetrized AND--OR search on two $16$-bit masks.
It evaluated $148{,}787$ distinct states and exported the policy used above.
The second program is the proof checker: its acceptance, together with the
finite reduction in Lemma~\ref{lem:certificate}, is the evidence for the
lower bound.  This separation between discovery and independently replayable
certification is also used in other exact finite computations; a recent
finite-geometry example reconstructs the relevant objects and constraints
before checking stored certificates \cite{ZengEtAl2026}.

As a further cross-check, a separate implementation reconstructed the
transversal masks and policy graph and also ran direct exact minimax.  The
root assertions that X can force a win by horizons $7,9,11$ were respectively
false, false and true.  This cross-check is consistent with the replayable
certificate but is not needed for the proof of the lower bound.

The certificate makes no assertion about boards larger than $4\times4$.
Moreover, the result is an exact finite computation rather than a structural
description of an optimal O strategy.  Finding a compact human-readable
defense, or extending the lower bound to arbitrary $n$, remains outside the
present result.

\paragraph{Data and computational disclosure.}
The policy, generator, independent verifier, verification reports and run
instructions accompany this article.  All decisive state and win tests use
integer arithmetic; no numerical tolerance or random seed is involved.
OpenAI Codex assisted with literature search, exact-search and verifier code
development, adversarial audit organization, and manuscript preparation.
Language-model output is not used as mathematical evidence: the lower bound
rests on the explicit finite reduction, the exported policy and its exact
verifier.  No proof assistant was used for the result reported here.

\bibliographystyle{abbrv}
\bibliography{references}

\end{document}
